% Part: normal-modal-logic
% Chapter: axioms-systems
% Section: logics-proofs

\documentclass[../../../include/open-logic-section]{subfiles}

\begin{document}

\olfileid{nml}{prf}{prf}

\olsection{\usetoken{P}{derivation} والنظم الجهية}

نضبط أولًا معنى !!a{derivation} في المنطق الجهوي النظامي.
فـ!!a{derivation}، على وجه الإجمال، متتالية من !!{formula}s:
كل !!{element} فيها إما حالة من \emph{بديهية}، وإما نتيجة
لـ!!{element}s سابقة بقاعدة من قواعد قليلة.
وبديهيات المنطق النظامي جميع أمثلة القضايا الصادقة صدقًا تحليليًا
\iftag{prvDiamond}{و\Ax{K} و~\Dual{}}{و~\Ax{K}}.
ومنها يتولد~$\Log{K}$، وهو أصغر منطق جهوي نظامي. فإن أردنا
نسقًا آخر زدنا بديهيات أخرى؛ وأما القواعد فلا تتبدل:
القياس الاستثنائي~\MP{} والإلزام~\Nec.

\begin{defn}
  إذا أعطي نسق جهوي $\Log{K} !A_\OLMathNumeral{1} \dots !A_\OLId{latin-ordinary}{n}$ و!!a{formula}~$!B$،
  قلنا إن $!B$ \emph{!!{derivable}} في $\Log{K} !A_\OLMathNumeral{1} \dots
  !A_\OLId{latin-ordinary}{n}$، وكتبنا $\Log{K} !A_\OLMathNumeral{1} \dots !A_\OLId{latin-ordinary}{n} \Proves !B$، إذا وفقط إذا
  وجدت !!{formula}s $!C_\OLMathNumeral{1}$، \dots، $!C_\OLId{latin-ordinary}{k}$ بحيث $!C_\OLId{latin-ordinary}{k} = !B$، وكان كل
  $!C_\OLId{latin-ordinary}{i}$ إما مثالا لقضية صادقة صدقا تحليليا، أو مثالا لواحدة من
  $\Ax{K}$،\iftag{prvDiamond}{ $\Dual$،}{} $!A_\OLMathNumeral{1}$، \dots، $!A_\OLId{latin-ordinary}{n}$،
  وإما لازما عن !!{formula}s سابقة بواسطة القاعدتين \MP{} أو~\Nec.
\end{defn}

وبالقضية الآتية يصير إثبات~$!B \in \Sigma$ ممكنًا بإقامة
$\Sigma$-!!{derivation} لـ~$!B$.

\begin{prop}
  $\Log{K} !A_\OLMathNumeral{1} \dots !A_\OLId{latin-ordinary}{n} = \Setabs{!B}{\Log{K} !A_\OLMathNumeral{1} \dots !A_\OLId{latin-ordinary}{n} \Proves !B}$.
\end{prop}

\begin{proof}
  نستعمل الاستقراء على طول !!{derivation}s لنبين أن
  $\Setabs{!B}{\Log{K} !A_\OLMathNumeral{1} \dots !A_\OLId{latin-ordinary}{n} \Proves !B} \subseteq
  \Log{K} !A_\OLMathNumeral{1} \dots !A_\OLId{latin-ordinary}{n}$.

  إذا كان طول !!{derivation} لـ~$!B$ هو~$\OLMathNumeral{1}$، فإنه يشتمل على
  !!{formula} واحدة. ولا يمكن أن تكون تلك !!{formula} لازمة عن صيغ
  سابقة بواسطة \MP{} أو \Nec، فلا بد أن تكون مثالا لقضية صادقة صدقا
  تحليليا، أو مثالا لـ\Ax{K}،\iftag{prvDiamond}{ $\Dual$،}{} أو مثالا
  لواحدة من $!A_\OLMathNumeral{1}$، \dots،~$!A_\OLId{latin-ordinary}{n}$. لكن $\Log{K}!A_\OLMathNumeral{1}\dots!A_\OLId{latin-ordinary}{n}$ يشتمل
  على هذه أيضا، ومن ثم $!B \in \Log{K}!A_\OLMathNumeral{1} \dots !A_\OLId{latin-ordinary}{n}$.

  وأما إذا كان طول !!{derivation} لـ$!B$ أكبر من~$\OLMathNumeral{1}$،
  فقد تكون~$!B$ بديهية كما سبق، أو ناتجة بواسطة~\MP{} أو~\Nec{}
  من !!{formula}s أسبق من آخر سطر في !!{derivation}.
  فإن كانت~$!B$ ناتجة من~$!C$ و$!C \lif !B$ بالقاعدة~\MP،
  كان~$!C$ و$!C \lif !B \in
  \Log{K}!A_\OLMathNumeral{1} \dots !A_\OLId{latin-ordinary}{n}$ بفرض الاستقراء. ومن الانغلاق تحت
  القياس الاستثنائي نحصل على~$!B \in \Log{K}!A_\OLMathNumeral{1} \dots !A_\OLId{latin-ordinary}{n}$.
  وإن كانت~$!B \ident \Box !C$ ناتجة من~$!C$ بالقاعدة~\Nec،
  كان~$!C
  \in \Log{K}!A_\OLMathNumeral{1} \dots !A_\OLId{latin-ordinary}{n}$ بفرض الاستقراء. والانغلاق
  تحت~$\Nec$ يعطي~$!B \in
  \Log{K}!A_\OLMathNumeral{1}\dots!A_\OLId{latin-ordinary}{n}$.

  وينتج الاحتواء العكسي من بيان أن
  $\Sigma = \Setabs{!B}{\Log{K} !A_\OLMathNumeral{1} \dots !A_\OLId{latin-ordinary}{n} \Proves !B }$ منطق جهوي
  نظامي يشتمل على جميع أمثلة $!A_\OLMathNumeral{1}$، \dots، $!A_\OLId{latin-ordinary}{n}$، ومن ملاحظة أن
  $\Log{K} !A_\OLMathNumeral{1} \dots !A_\OLId{latin-ordinary}{n}$ هو بالتعريف أصغر منطق من هذا القبيل.
  \begin{enumerate}
    \item كل قضية صادقة صدقا تحليليا~$!B$ مثال لمثل هذه القضية، ولذلك
      $\Log{K}!A_\OLMathNumeral{1}\dots!A_\OLId{latin-ordinary}{n} \Proves !B$، فتشتمل $\Sigma$ على جميع
      القضايا الصادقة صدقا تحليليا.
    \item إذا كان $\Log{K}!A_\OLMathNumeral{1}\dots!A_\OLId{latin-ordinary}{n} \Proves !C$ و
      $\Log{K}!A_\OLMathNumeral{1}\dots!A_\OLId{latin-ordinary}{n} \Proves !C \lif !B$، فعندئذ
      $\Log{K}!A_\OLMathNumeral{1}\dots!A_\OLId{latin-ordinary}{n} \Proves !B$: اجمع !!{derivation} لـ
      $!C$ مع اشتقاق $!C \lif !B$، وأضف السطر~$!B$. ويسوغ السطر
      الأخير بـ\MP{}. وهكذا تكون $\Sigma$ مغلقة تحت القياس الاستثنائي.
    \item إذا كان لـ$!B$ !!a{derivation}، فإن لكل مثال إحلال لـ~$!B$
      أيضا !!a{derivation}: طبّق الإحلال على كل !!{formula} في
      !!{derivation}. (تمرين: أثبت بالاستقراء على طول !!{derivation}s
      أن الناتج أيضا !!{derivation} صحيح.) ومن ثم تكون $\Sigma$
      مغلقة تحت الإحلال المنتظم. (وقد أثبتنا الآن أن $\Sigma$ تحقق
      جميع شروط المنطق الجهوي.)
    \item لدينا $\Log{K}!A_\OLMathNumeral{1}\dots!A_\OLId{latin-ordinary}{n} \Proves \Ax{K}$، ومن ثم
      $\Ax{K} \in \Sigma$.
    \tagitem{prvDiamond}{لدينا
      $\Log{K}!A_\OLMathNumeral{1}\dots!A_\OLId{latin-ordinary}{n} \Proves \Dual$، ومن ثم $\Dual \in \Sigma$.}{}
    \item إذا كان $\Log{K}!A_\OLMathNumeral{1}\dots!A_\OLId{latin-ordinary}{n} \Proves !C$، كان السطر
      الإضافي~$\Box !C$ مسوغا بـ~\Nec. ولذلك تكون $\Sigma$ مغلقة
      تحت~\Nec، ومن ثم تكون $\Sigma$ نظامية.
  \end{enumerate}
\end{proof}

\end{document}
